% ═══════════════════════════════════════════════════════════════
% ARBEITSBLATT: Numbers Einstieg
% Informatik Klasse 7 | Stunde 4
% ═══════════════════════════════════════════════════════════════

\documentclass[11pt, a4paper]{article}

\usepackage[utf8]{inputenc}
\usepackage[T1]{fontenc}
\usepackage[ngerman]{babel}

\usepackage{helvet}
\renewcommand{\familydefault}{\sfdefault}

\usepackage[margin=1.5cm, top=1.8cm, bottom=1.5cm]{geometry}
\usepackage{enumitem}
\usepackage{tabularx}
\usepackage{booktabs}
\usepackage{array}
\usepackage{tikz}
\usepackage{fancyhdr}
\usepackage{xcolor}
\usepackage{tcolorbox}
\tcbuselibrary{skins}
\usepackage{amssymb}
\usepackage{qrcode}
\usepackage{hyperref}
\hypersetup{colorlinks=false, pdfborder={0 0 0}}
\usepackage{multicol}

% ═══════════════════════════════════════════════════════════════
% FARBEN
% ═══════════════════════════════════════════════════════════════

\definecolor{informatik}{HTML}{b35c00}
\definecolor{hellgrau}{HTML}{F5F5F5}
\definecolor{mittelgrau}{HTML}{888888}

\colorlet{fachfarbe}{informatik}

% ═══════════════════════════════════════════════════════════════
% KOPFZEILE
% ═══════════════════════════════════════════════════════════════

\pagestyle{fancy}
\fancyhf{}
\renewcommand{\headrulewidth}{0.4pt}
\renewcommand{\footrulewidth}{0pt}
\lhead{\small Informatik Klasse 7 | Numbers Einstieg}
\rhead{\small Name: \underline{\hspace{4cm}}}

% ═══════════════════════════════════════════════════════════════
% BOXEN
% ═══════════════════════════════════════════════════════════════

\newtcolorbox{merkkasten}[1][]{
    colback=hellgrau,
    colframe=hellgrau,
    boxrule=0pt,
    arc=0pt,
    left=8pt, right=8pt, top=6pt, bottom=5pt,
    borderline north={2pt}{0pt}{fachfarbe},
    before skip=6pt, after skip=6pt,
    #1
}

\newtcolorbox{lueckenkasten}[1][]{
    colback=white,
    colframe=fachfarbe,
    boxrule=1pt,
    arc=0pt,
    left=6pt, right=6pt, top=6pt, bottom=6pt,
    before skip=4pt, after skip=4pt,
    #1
}

% Einsetzstrich
\newcommand{\luecke}[1]{\underline{\hspace{#1}}}

% URL
\newcommand{\lpurl}{https://devbox42.github.io/sciencesim/informatik/kl07/tabellenkalkulation/04-numbers-einstieg/LP-04-numbers-einstieg.html}

\setlength{\parindent}{0pt}
\setlength{\parskip}{3pt}
\setlength{\columnsep}{0.8cm}

% ═══════════════════════════════════════════════════════════════
% DOKUMENT
% ═══════════════════════════════════════════════════════════════

\begin{document}

% Header mit QR-Code
\begin{minipage}[t]{0.78\textwidth}
    \vspace{0pt}
    {\Large\bfseries\textcolor{fachfarbe}{Numbers Einstieg – Erste Schritte mit Apple Numbers}}\\[0.3em]
    {\small Arbeitsblatt | Zum Abheften}
\end{minipage}%
\hfill%
\begin{minipage}[t]{0.18\textwidth}
    \vspace{0pt}
    \raggedleft
    \href{\lpurl}{\fcolorbox{fachfarbe}{white}{\qrcode[height=18mm]{\lpurl}}}
\end{minipage}

\vspace{0.2em}
\hrule
\vspace{0.4em}

\begin{multicols}{2}

% ═══════════════════════════════════════════════════════════════
% KASTEN 1: Programme starten mit Spotlight
% ═══════════════════════════════════════════════════════════════

\textbf{\textcolor{fachfarbe}{Kasten 1: Programme mit Spotlight starten}}

\begin{lueckenkasten}
\textbf{Spotlight-Suche} ist die schnellste Art, Programme auf dem Mac zu öffnen.

\vspace{0.2em}
Tastenkombination: \luecke{1.5cm} + \luecke{2cm}

\vspace{0.2em}
Dann tippen: \luecke{2.5cm} und \textbf{Enter} drücken
\end{lueckenkasten}

% ═══════════════════════════════════════════════════════════════
% KASTEN 2: App-Wechsel mit ⌘+Tab
% ═══════════════════════════════════════════════════════════════

\textbf{\textcolor{fachfarbe}{Kasten 2: Zwischen Apps wechseln (Cmd+Tab)}}

\begin{lueckenkasten}
Der \textbf{App-Umschalter} zeigt alle offenen Programme als Symbole an.

\vspace{0.2em}
Tastenkombination: \luecke{1.5cm} + \luecke{1.5cm}

\vspace{0.2em}
\textit{Tipp:} \luecke{1.5cm} gedrückt halten, mit \luecke{1.5cm} durchschalten.
\end{lueckenkasten}

% ═══════════════════════════════════════════════════════════════
% KASTEN 3: Formeln in Numbers eingeben
% ═══════════════════════════════════════════════════════════════

\textbf{\textcolor{fachfarbe}{Kasten 3: Formeln eingeben (Gesamtpreis)}}

\begin{lueckenkasten}
Formeln beginnen \textbf{immer} mit dem Zeichen: \luecke{1cm}

\vspace{0.2em}
\textbf{Beispiel:} Preis (B2) $\times$ Menge (C2):

\texttt{=}\luecke{1cm}\texttt{*}\luecke{1cm}

\vspace{0.2em}
Nach Eingabe: \luecke{2cm} drücken zur Bestätigung
\end{lueckenkasten}

% ═══════════════════════════════════════════════════════════════
% KASTEN 4: Formeln kopieren (Aufziehen)
% ═══════════════════════════════════════════════════════════════

\textbf{\textcolor{fachfarbe}{Kasten 4: Formeln kopieren (Aufziehen)}}

\begin{lueckenkasten}
\textbf{Methode 1 – Aufziehen:}

Den gelben \luecke{2cm} nach unten ziehen.

\vspace{0.2em}
\textbf{Methode 2 – Tastatur:}

Kopieren: \luecke{1.5cm} + \luecke{1cm}

Einfügen: \luecke{1.5cm} + \luecke{1cm}
\end{lueckenkasten}

\columnbreak

% ═══════════════════════════════════════════════════════════════
% KASTEN 5: SUMME-Funktion für Gesamtsumme
% ═══════════════════════════════════════════════════════════════

\textbf{\textcolor{fachfarbe}{Kasten 5: SUMME-Funktion (Gesamtsumme)}}

\begin{lueckenkasten}
\texttt{=SUMME(D2:D5)} addiert alle Werte von D2 bis D5.

\vspace{0.2em}
Der \textbf{Doppelpunkt} \texttt{:} bedeutet: \luecke{2cm}

\vspace{0.2em}
\textbf{Ergebnis:} Die Summe aller \luecke{3cm}
\end{lueckenkasten}

% ═══════════════════════════════════════════════════════════════
% KASTEN 6: Speichern mit ⌘+S
% ═══════════════════════════════════════════════════════════════

\textbf{\textcolor{fachfarbe}{Kasten 6: Datei speichern (Cmd+S)}}

\begin{lueckenkasten}
\textbf{Speichern-Tastenkombination:}

\luecke{1.5cm} + \luecke{1cm}

\vspace{0.2em}
\textbf{Dateiname:} Kiosk\_Vorname

\textbf{Speicherort:} \luecke{3cm}

\textbf{Dateiendung:} \texttt{.}\luecke{2cm}
\end{lueckenkasten}

% ═══════════════════════════════════════════════════════════════
% KASTEN 7: Diagramm erstellen (Symbolleiste)
% ═══════════════════════════════════════════════════════════════

\textbf{\textcolor{fachfarbe}{Kasten 7: Säulendiagramm erstellen}}

\begin{lueckenkasten}
\textbf{1.} Daten (Produkt + Gesamtpreis) \luecke{2.5cm}

\textbf{2.} In Symbolleiste auf \luecke{2.5cm} klicken

\textbf{3.} Kategorie \luecke{2cm}, dann \luecke{2.5cm} wählen
\end{lueckenkasten}

% ═══════════════════════════════════════════════════════════════
% KASTEN 8: Abgabe in der Bildungscloud
% ═══════════════════════════════════════════════════════════════

\textbf{\textcolor{fachfarbe}{Kasten 8: Abgabe in der Bildungscloud}}

\begin{lueckenkasten}
\textbf{1.} Im \textbf{Finder} die \texttt{.numbers}-Datei finden

\textbf{2.} In der Bildungscloud: \luecke{2.5cm} öffnen

\textbf{3.} Auf \luecke{3cm} klicken

\textbf{4.} Datei hochladen, dann \luecke{2.5cm}
\end{lueckenkasten}

\end{multicols}

\vspace{0.3em}
\hrule
\vspace{0.4em}

% ═══════════════════════════════════════════════════════════════
% ZUSAMMENFASSUNG
% ═══════════════════════════════════════════════════════════════

\begin{merkkasten}
\textbf{Zusammenfassung:} Mit \textbf{Cmd+Leertaste} öffnest du Spotlight und startest Numbers. Mit \textbf{Cmd+Tab} wechselst du zwischen Browser und Numbers. \textbf{Formeln} beginnen mit \texttt{=}. Mit dem \textbf{gelben Punkt} ziehst du Formeln nach unten. \textbf{Cmd+S} speichert deine Datei. Die fertige \texttt{.numbers}-Datei lädst du in der \textbf{Bildungscloud} hoch.
\end{merkkasten}

\end{document}
